\section{The source and its observation channels}\label{sec:observation}
For the equilibrium example, the relevant symmetry concerns the particle before its signals pass through the sensors. Covariance measures how two fluctuations vary together, including their scale. Let $C^X_{ij}(k)$ be the covariance of source channels $i$ and $j$ at a time separation, or lag, $k$. Because means and covariances determine the Gaussian law, their symmetry can establish reversibility. After subtraction of the mean, a stationary Gaussian process whose coordinates have a common time-reversal parity satisfies
\[
 \text{reversibility}\quad\Longleftrightarrow\quad
 C^X_{ij}(k)=C^X_{ij}(-k).
\]
The equality must hold for every channel pair and lag. It makes the full Gaussian law invariant under reversal, since the mean and covariance determine that law. Equivalently, the source cross-spectral measure is real. The common-parity condition includes position coordinates, which keep their signs under reversal, and separately velocity coordinates, which change their signs. A mixture of positions and velocities requires parity signs in the covariance criterion and is outside the convention used here.

The spectrum expresses the same covariance relations by frequency, where a detector timing change becomes a phase factor. A linear sensor records a weighted sum of source values. This operation changes both the amplitude and timing of the signal. Write its transfer function as $H_i(\omega)$ for channel $i$ and frequency $\omega$. A pure delay has response $e^{-id_i\omega}$, while the echo in Fig.~\ref{fig:sensor-example} has response $1+0.03e^{-i\omega}$. If $\mu$ is the matrix spectral measure of the source, the recorded spectral measure is
\begin{equation}
 F=D_H\mu D_H^*+F_E,\qquad D_H=\operatorname{diag}(H_i).
 \label{eq:main-model}
\end{equation}
The star denotes conjugate transpose. The source and errors are jointly Gaussian and stationary. Errors are uncorrelated with the source and between channels at every lag, so their spectral measure $F_E$ is diagonal. Unknown constant means are allowed. For $i\ne j$, Eq.~\eqref{eq:main-model} multiplies the source cross spectrum by $H_i\overline H_j$. Its phase can therefore change even when the source spectrum is real.

We use $C_{ij}(k)=\operatorname{Cov}\{Y_i(t+k),Y_j(t)\}=\int e^{ik\omega}F_{ij}(d\omega)$ for the recorded covariance. Discrete-time densities use $d\omega/(2\pi)$ on $[-\pi,\pi]$. Spectral measures also allow contributions concentrated at individual frequencies. Equation~\eqref{eq:main-model} describes filters acting on the sampled source sequences. Filtering in continuous time before sampling does not always give this factorization. Distinct continuous frequencies can become one sampled frequency while retaining different detector phases. This aliasing will matter both for the ambiguity result and for the interpretation of phase cancellation.

\subsection{Equilibrium and driving in a trapped particle}
We use a Brownian gyrator to connect the statistical symmetry to a physical source~\cite{filliger2007,cerasolispectra2022,argun2017}. A particle moves in a two-dimensional harmonic trap. The restoring force couples its coordinates, and each coordinate receives random forcing from a heat bath. In the overdamped limit, where inertia is neglected, the dimensionless dynamics are
\begin{equation}
 \begin{gathered}
 dX=-AX\,dt+\sqrt{2D}\,dW_t,\\
 A=\begin{pmatrix}1&1/2\\1/2&1\end{pmatrix},\quad
 D=\operatorname{diag}(T_1,T_2).
 \end{gathered}
 \label{eq:ou-model}
\end{equation}
Here $dW_t$ supplies the random forcing, and the stationary covariance satisfies $A\Sigma+\Sigma A^{\mathsf T}=2D$. Equal bath temperatures give equilibrium. Unequal temperatures produce a circulating stationary probability current: trajectories favor one direction around the trap although the position distribution stays constant. We sample at interval one and then apply $H_1=1$, $H_2=1+0.03e^{-i\omega}$. The equilibrium curve in Fig.~\ref{fig:sensor-example} therefore isolates the effect of unequal detectors within Eq.~\eqref{eq:main-model}. The same source model with unequal temperatures will test whether detector information permits detection of physical driving.

Throughout the paper, the null hypothesis is a reversible source observed through the allowed channels. A test must control the probability of rejecting this explanation whenever it is true. Failure to reject does not establish reversibility. We also distinguish nonrejection from abstention: a test abstains when a required bound or numerical comparison is unavailable. A source outside the null class is an alternative. Its rejection probability is the detection probability, or power. All records, including abstentions, remain in the denominators of the reported rejection fractions. Table~\ref{tab:notation} in Methods collects the main symbols.

\section{What a finite recording cannot decide}\label{sec:ambiguity}
The equilibrium example suggests a possible remedy: collect enough data to resolve the source and detector contributions separately. Without restrictions on the source or detectors, this remedy fails. A finite record contains only finitely many covariance lags. A delayed reversible source can agree with an asymmetric source at every retained lag and differ only outside the record.

Consider two explanations for the same two-time record. In the first, the source is asymmetric and the sensors add no delay. In the second, the source is reversible and one sensor delays it. Each channel separately has variance one and no covariance between distinct times, called a unit white marginal spectrum. The first explanation has cross spectrum
\[
 F^A_{12}=\tfrac1{10}+\tfrac15e^{-i\omega}.
\]
A reversible source with $S^R_{12}=\tfrac15+\tfrac15\cos\omega$, observed after delaying channel 1 by one sample, instead gives
\[
 F^R_{12}=e^{-i\omega}S^R_{12}
 =\tfrac1{10}+\tfrac15e^{-i\omega}+\tfrac1{10}e^{-2i\omega}.
\]
The spectral eigenvalues are at least $7/10$ and $3/5$, respectively, so both constructions are positive and define valid Gaussian processes. Their cross-covariances are
\[
\begin{array}{c|ccccc}
 k&-2&-1&0&1&2\\\hline
 C^A_{12}(k)&0&0&1/10&1/5&0\\
 C^R_{12}(k)&0&0&1/10&1/5&1/10
\end{array}
\]
A two-time record sees only lags $-1,0,1$, the middle three columns. Its complete covariance, and hence its Gaussian law, agrees under both explanations. The extra term in the second spectrum changes only the last column shown. Lag two would distinguish the explanations, but is absent from that record.


The following result extends that construction from a two-time example to every finite stationary Gaussian law with positive-definite covariance. Positive definiteness excludes an exact linear constraint between recorded fluctuations. The constructed sources have finite dependence: their covariances vanish beyond a finite lag.
\begin{theorem}[Finite-law ambiguity]\label{thm:main-finite-law}
Fix a positive record length and a nonempty finite collection of stationary Gaussian record laws with positive-definite covariance matrices and the same number of channels. One vector of nonnegative integer channel delays gives a reversible Gaussian source representation for every law in the collection. Each source has a strictly positive matrix trigonometric polynomial spectral density and finite dependence. The sources preserve the channel means and variances. No observation error is needed.
\end{theorem}
Thus even the exact finite joint distribution cannot exclude a reversible source in this unrestricted class. A test with false-rejection probability at most $\alpha$ throughout that class has detection probability at most $\alpha$ against every nonsingular finite stationary Gaussian alternative. The proof constructs a positive reversible spectrum, then corrects the retained moments exactly (Supplemental Material, abbreviated SM, Sec.~S1, Theorem~S1 and Proposition~S1). An approximate covariance match would not suffice.

For two channels and $N\ge2$, generic retained covariances require integer relative delay at least $N-1$. The exceptional records satisfy $C_{12}(\Delta-1)=C_{12}(\Delta+1)$ for an integer $|\Delta|\le N-2$. The minimum scale is attainable under an additional covariance condition (SM Sec.~S1.2). The source and delays can depend on record length; the result concerns a fixed finite law, not one representation valid at all lengths.

Large delays are not the only source of ambiguity. Suitable prescribed real delays, whose nonreference relative delays are rationally independent with the sampling interval, can be arbitrarily small when continuous frequencies are unrestricted. This means that no nonzero integer combination of those delays and the interval is zero. Aliased frequencies then supply the required channel phases (SM Sec.~S2, Theorem~S3 and Corollary~S2). A spectral envelope instead bounds the effect of small skew (Proposition~S3). Bandwidth control therefore matters as well as timing accuracy~\cite{hansenSargent1983,parratobar,rutten}. The next two sections restrict the ambiguity through response calibration or phase cancellation, together with quantitative covariance bounds.
